\section{Experimental Setup}
\label{sec:experimental}

We design experiments to answer the following research questions:

\textbf{RQ1:} Does a non-monotonic dose-response relationship exist between perplexity filtering thresholds and benchmark performance?

\textbf{RQ2:} Does the noise dilution mechanism explain why intermediate thresholds outperform extremes?

\textbf{RQ3:} Where is the optimal perplexity threshold, and how precisely can it be identified?

\textbf{RQ4:} Do optimal thresholds transfer across model scales (125M to 1B)?

\textbf{RQ5:} Does CPDR-optimized configuration outperform industry-standard RedPajama defaults?

\subsection{Dataset}

We use RedPajama-v2 \citep{together2023redpajama}, a large-scale web corpus representative of modern LLM training data. The dataset provides raw text samples with quality annotations enabling controlled filtering experiments.

\begin{table}[h]
\centering
\caption{Dataset properties.}
\label{tab:dataset}
\begin{tabular}{@{}ll@{}}
\toprule
Property & Value \\
\midrule
Source & Common Crawl + curated sources \\
Total Documents & 30B+ \\
Our Sample & 10B tokens (target); 5M--10M tokens (PoC) \\
Language & English \\
Quality Signal & KenLM 5-gram perplexity \\
\bottomrule
\end{tabular}
\end{table}

\textbf{Why RedPajama:} The dataset is (1) representative of production LLM training data, (2) provides raw quality signals enabling threshold experimentation, and (3) has established baseline configurations in the literature for comparison.

\subsection{Baselines}

We compare against:

\textbf{No Filtering (p0):} Raw corpus without perplexity filtering. Included to establish lower bound and test noise dilution hypothesis.

\textbf{RedPajama Defaults:} Literature-standard perplexity and deduplication thresholds from RedPajama documentation. Represents current industry practice.

\textbf{CPDR-Optimized:} Curation parameters selected via our dose-response sweep methodology. Tests whether systematic optimization improves over heuristic choices.

\subsection{Implementation Details}

\textbf{Model Architecture:} GPT-2 variants (125M, 350M, 1B parameters) using HuggingFace Transformers implementation \citep{radford2019gpt2}.

\textbf{Training Configuration:}
\begin{itemize}
    \item Learning rate: $6 \times 10^{-4}$ with cosine decay
    \item Batch size: 512 sequences
    \item Warmup: 2000 steps
    \item Training tokens: 10B (target), 5M--10M (PoC validation)
    \item Random seeds: 42 (primary), 123, 456 (variance estimation)
\end{itemize}

\textbf{Compute Resources:} Training performed on $4 \times$ A100 GPUs. Single 125M configuration requires approximately 8 GPU-hours at PoC scale.

\textbf{Perplexity Scoring:} KenLM 5-gram model trained on Wikipedia following CCNet methodology \citep{wenzek2020ccnet}. Samples filtered by percentile threshold (p0 through p90).

\subsection{Evaluation Protocol}

\textbf{Benchmark Ensemble:} We evaluate on four reasoning benchmarks: HellaSwag, ARC-Easy, PIQA, and WinoGrande. Following scaling laws literature \citep{hoffmann2022chinchilla}, we compute a benchmark ensemble score as the first principal component (PC1) of accuracy scores, reducing noise from individual benchmark variance.

\begin{table}[h]
\centering
\caption{Evaluation benchmarks.}
\label{tab:benchmarks}
\begin{tabular}{@{}lll@{}}
\toprule
Benchmark & Task Type & Metric \\
\midrule
HellaSwag & Commonsense reasoning & Accuracy \\
ARC-Easy & Science QA & Accuracy \\
PIQA & Physical reasoning & Accuracy \\
WinoGrande & Coreference resolution & Accuracy \\
\bottomrule
\end{tabular}
\end{table}

\textbf{Contamination Verification:} We apply Min-K\%++ \citep{shi2024minkpp} membership inference to verify evaluation set integrity before reporting scores.

\textbf{Statistical Analysis:} Polynomial regression with AIC/BIC model selection for dose-response characterization. Bootstrap resampling (1000 iterations) for confidence intervals on optimal threshold estimates.
